%; whizzy paragraph -pdf xpdf -latex ./whizzypdfptex.sh
%; whizzy-paragraph "^\\\\begin{frame}"
% latex beamer presentation.
% platex, latex-beamer でコンパイルすることを想定。 

%     Tokyo Debian Meeting resources
%     Copyright (C) 2009 Junichi Uekawa
%     Copyright (C) 2009 Nobuhiro Iwamatsu

%     This program is free software; you can redistribute it and/or modify
%     it under the terms of the GNU General Public License as published by
%     the Free Software Foundation; either version 2 of the License, or
%     (at your option) any later version.

%     This program is distributed in the hope that it will be useful,
%     but WITHOUT ANY WARRANTY; without even the implied warranty of
%     MERCHANTABILITY or FITNESS FOR A PARTICULAR PURPOSE.  See the
%     GNU General Public License for more details.

%     You should have received a copy of the GNU General Public License
%     along with this program; if not, write to the Free Software
%     Foundation, Inc., 51 Franklin St, Fifth Floor, Boston, MA  02110-1301 USA

\documentclass[cjk,dvipdfmx,12pt]{beamer}
\usetheme{Tokyo}
\usepackage{monthlypresentation}

%  preview (shell-command (concat "evince " (replace-regexp-in-string "tex$" "pdf"(buffer-file-name)) "&"))
%  presentation (shell-command (concat "xpdf -fullscreen " (replace-regexp-in-string "tex$" "pdf"(buffer-file-name)) "&"))
%  presentation (shell-command (concat "evince " (replace-regexp-in-string "tex$" "pdf"(buffer-file-name)) "&"))

%http://www.naney.org/diki/dk/hyperref.html
%日本語EUC系環境の時
\AtBeginDvi{\special{pdf:tounicode EUC-UCS2}}
%シフトJIS系環境の時
%\AtBeginDvi{\special{pdf:tounicode 90ms-RKSJ-UCS2}}

\title{lxc を使ってみる。}
\subtitle{今年最後のネタです。}
\author{まえだこうへい mkouhei@debian.or.jp \\IRC nick: mkouhei}
\date{2009年12月12日}
\logo{\includegraphics[width=8cm]{image200607/openlogo-light.eps}}

\begin{document}

\frame{\titlepage{}}

\begin{frame}{lxc とは何ぞや？}
\begin{itemize}
 \item 正式名称 Linux Containers
 \item lxc の構成
       \begin{itemize}
	\item コンテナ自体が稼働するカーネルの機能
	\item lxc を管理するための lxc-* コマンドラインツール
       \end{itemize}
 \item Daniel Lezcano 氏が開発、メンテナンス。GPL2で配布。
 \item Linux Kernel 2.6.29 で lxc に必要なパッチが mainline にマージされ
       た。
\end{itemize}
\end{frame}

\begin{frame}{他との比較}
\begin{itemize}
 \item OpenVZ, Linux-VServer
 \item Solaris Containers
 \item FreeBSD jail
 \item lxc で提供するサービス
       \begin{itemize}
	\item システムコンテナ：いわゆるOSの仮想化
	\item アプリケーションコンテナ：chroot でのプロセスの分離
       \end{itemize}
\end{itemize}
\end{frame}

\begin{frame}{lxc を使ってみた動機。}
\begin{itemize}
 \item 連載も終わったし、KVM 以外も見てみるか。
 \item あまり使われてないマイナーなのもの。
 \item lguest だと 64bit 環境で使えない。
 \item たまたま dW で見つけた。
 \item でもある程度まとまったのはまだ dW の記事しかない。
 \item 何より手軽！
\end{itemize}
\end{frame}

\emtext{導入する}
\begin{frame}[containsverbatim]{パッケージインストール}
\begin{commandline}
$ sudo apt-get install lxc iproute debootstrap
\end{commandline}
\end{frame}

\begin{frame}[containsverbatim]{確認する}

\begin{commandline}
$ lxc-checkconfig
--- Namespaces ---
Namespaces: enabled
Utsname namespace: enabled
Ipc namespace: enabled
Pid namespace: enabled
User namespace: enabled
Network namespace: enabled
Multiple /dev/pts instances: enabled

--- Control groups ---
Cgroup: enabled
Cgroup namespace: enabled
Cgroup device: enabled
Cgroup sched: enabled
Cgroup cpu account: enabled
Cgroup memory controller: enabled
Cgroup cpuset: enabled

--- Misc ---
Veth pair device: enabled
Macvlan: enabled
File capabilities: enabled
\end{commandline}

なんか足りなければ、カーネルリビルドしとけ。
\end{frame}

\begin{frame}[containsverbatim]{３分もいらないぜ。}
即席コンテナで稼働確認する。

\begin{commandline}
$ sudo lxc-execute -n hoge -f \
 /usr/share/doc/lxc/examples/lxc/lxc-macvlan.conf /bin/bash
$ sudo lxc-info -n hoge
'hoge' is RUNNING
$ sudo lxc-stop -n hoge <-別のシェルで実行。
\end{commandline}
\end{frame}

\begin{frame}[containsverbatim]{システムコンテナを使う。}
簡単に Debian 環境を増やせて、簡単に起動できるのが良い。

やること。
\begin{itemize}
 \item ネットワークの設定
 \item cgroupファイルシステムのマウント
 \item コンテナの作成
 \item コンテナのファイルの微調整
 \item コンテナの起動
\end{itemize}
\end{frame}

\emtext{lxc の仕組みの話}
\begin{frame}{lxc のライフサイクル}

他の仮想マシンのライフサイクルと大差ない。

\includegraphics[width=1.0\hsize]{image200912/lxc/lxc-lifecycle.eps}

\end{frame}

\begin{frame}{lxc のライフサイクル}

{\footnotesize
\begin{tabular}[t]{|p{5em}|p{5em}|p{6em}|p{11em}|}
\hline
status &コマンド &実行後の状態 &備考 \\
\hline
\hline
コンテナの起動 &lxc-start &STARTING, RUNNING &システムコンテナ \\
&lxc-execute & &アプリケーションコンテナ\\
\hline
コンテナの一時停止 &lxc-freeze &FROZEN & \\
\hline
コンテナの再開 &lxc-unfreeze &RUNNING & \\
\hline
コンテナの停止 &lxc-stop &STOPPING, STOPPED & \\
\hline
コンテナの再起動 &lxc-restart &STOPPING, STOPPED & \\
& &STARTING, RUNNING & \\
\hline
コンテナの作成 &lxc-create &STOPPED &lxc-debian は内部でlxc-create を実
             行。 \\
\hline
コンテナの破棄 &lxc-destroy & & \\
\hline
\end{tabular}}

\end{frame}

\begin{frame}[containsverbatim]{リソースの管理}
\begin{itemize}
\item cgroup(Control Group)で行う。
\item インタフェースはlxc-cgroup コマンド
\item マウントした cgroup ファイルシステムからアクセスできる。
\begin{commandline}
cgroup  /var/local/cgroup  cgroup  defaults  0  0
\end{commandline}
\item システムコンテナの rootfs を置くのはどこでも良い。
\item Read, Write の権限は この/var/local/cgroup 以下のファイルのアクセ
      ス権で決まる。
\end{itemize}
\end{frame}

\emtext{おまけ}
\begin{frame}{libvirt の話}
\begin{itemize}
\item lxc は libvirt のサポート対象に既になっている。\\
      んだが、うまくいかない。
\item システムコンテナを定義するとこまではできた。
\end{itemize}
\end{frame}


\begin{frame}[containsverbatim]{libvirt の話}
\begin{commandline}
(test.xml)
<domain type='lxc'>
  <name>test</name>
  <memory>500000</memory>
  <os>
     <type>exe</type>
     <init>/sbin/init</init>
  </os>
  <vcpu>1</vcpu>
  <clock offset='utc'/>
  <on_poweroff>destroy</on_poweroff>
  <on_reboot>restart</on_reboot>
  <on_crash>destroy</on_crash>
  <devices>
     <emulator>/usr/lib/libvirt/libvirt_lxc</emulator>
     <filesystem type='mount'>
       <source dir='/var/cache/lxc/debian/rootfs.test'/>
       <target dir='/'/>
     </filesystem>
     <interface type='bridge'>
       <source bridge='br0'/>
       <target dev='veth_99999'/>
       <mac address='00:11:22:33:44:55'/>
     </interface>
     <console type='pty' />
  </devices>
</domain>
\end{commandline}
\end{frame}

\begin{frame}[containsverbatim]{libvirt の話}
\begin{commandline}
$ sudo virsh --connect lxc:/// define text.xml
$ sudo virsh --connect lxc:/// start lxc
$ sudo virsh --connect lxc:/// list
 Id Name                 State
----------------------------------
  3 test                 no state
\end{commandline}
なんでやねん。
\end{frame}

\begin{frame}[containsverbatim]{lxc-*コマンドでの運用を考える。}
\begin{itemize}
\item lxc-* のコマンドは、1コンテナに対するアクセスしかできない。
\item 全体を把握するには不便。
\item lxc-ls コマンドあたりを活用すると便利。
\end{itemize}
\end{frame}

\begin{frame}[containsverbatim]{lxc-*コマンドでの運用を考える。}
\begin{commandline}
#!/bin/sh

for i in `lxc-ls`
do
if (lxc-info -n $i | grep -q "STOPPED" ) then
   echo -en $i"\t:" 
   lxc-start -n $i -d
   test $? -eq 0 && echo starting ok. 
   sleep 3
fi
done

exit 0
\end{commandline}
\end{frame}

\begin{frame}{まとめ}
\begin{itemize}
 \item lxc とは コンテナ。
 \item lxc を導入してみる。ほとんど手間いらず。
 \item lxc の仕組み。cgroupの具体的なサンプル。
 \item libvirt はもうちょい先かな。
\end{itemize}

\end{frame}


\end{document}

;;; Local Variables: ***
;;; outline-regexp: "\\([ 	]*\\\\\\(documentstyle\\|documentclass\\|emtext\\|section\\|begin{frame}\\)\\*?[ 	]*[[{]\\|[]+\\)" ***
;;; End: ***
